\chapter{KẾT LUẬN VÀ HƯỚNG PHÁT TRIỂN}
Nghiên cứu này đã thực hiện quy trình đánh giá thực nghiệm toàn diện và độc lập đối với các mô hình học sâu tiên tiến cho hai bài toán cốt lõi: Phân vùng cấu trúc và phát hiện bất thường trên ảnh MRI cột sống thắt lưng. Kết quả thực nghiệm trên các tập dữ liệu chuẩn hóa đã cung cấp cơ sở định lượng vững chắc cho việc lựa chọn giải pháp kỹ thuật phù hợp với mục tiêu xây dựng hệ thống hỗ trợ gán nhãn tự động.
\section{Tổng kết}
Dựa trên các chỉ số đánh giá và phân tích định lượng, nghiên cứu rút ra các kết luận chính sau: \\

\textbf{Đối với bài toán phân vùng cột sống:} Mô hình \textbf{TotalSpineSeg} được lựa chọn sau khi so sánh với 3 mô hình khác (SPINEPS, MedCLIP-SAMv2, SpineNetV2).
    \begin{itemize}
        \item Mô hình đạt độ chính xác dương tính (Precision) lên tới \textbf{97.87\%}, đảm bảo vùng được trích xuất gần như không nhiễu, vượt trội so với SPINEPS (52.17\%), MedCLIP-SAMv2 (29.16\%) và SpineNetV2 (51.76\%).
        \item Mặc dù chỉ số Recall thấp hơn (56.49\%) do khác biệt về định nghĩa nhãn thực tế (TotalSpineSeg chỉ tập trung thân đốt sống, không tính gai sau), nhưng độ ổn định và độ chính xác bề mặt (ASSD = 6.85 mm thấp nhất) khiến đây là lựa chọn tối ưu để trích xuất ROI sạch cho các bước xử lý sau.
    \end{itemize}

\textbf{Đối với bài toán phát hiện bất thường (chọn mô hình nền):} Mô hình \textbf{SpineNetV2} được lựa chọn làm backbone cho mô-đun grading.
    \begin{itemize}
        \item Trên 3 nhóm bệnh lý chính (Hẹp ống sống, Hẹp lỗ liên hợp trái và phải), SpineNetV2 đạt F1-Score trung bình \textbf{77.50\%} --- thấp hơn Ning Shen (SOTA, F1 80.37\%) khoảng 3\% nhưng được ưu tiên nhờ mã nguồn mở, kiến trúc đơn giản, dễ tinh chỉnh sâu. Ning Shen là ensemble nhiều mô hình lớn, khó duy trì.
        \item MedGemma bị loại do Recall quá thấp (6--21\%) trên các bệnh lý cột sống thắt lưng do chưa được chuyên biệt hóa.
        \item \emph{Lưu ý:} thí nghiệm chọn mô hình chỉ đánh giá 3/5 conditions vì SpineNetV2 nguyên bản không hỗ trợ Left/Right Subarticular Stenosis --- hạn chế này đã được khắc phục trong đóng góp tinh chỉnh ở phần dưới.
    \end{itemize}

\textbf{Đối với đóng góp tinh chỉnh sâu mô-đun grading:} Kiến trúc Hybrid CBAM + BiomedCLIP được đề xuất và đánh giá kĩ trên hai dataset độc lập.
    \begin{itemize}
        \item \textbf{Kiến trúc hỗ trợ 5 conditions:} thiết kế mới gồm 5 classification head, bổ sung Left/Right Subarticular Stenosis mà SpineNetV2 nguyên bản còn thiếu --- đạt độ phủ đầy đủ chuẩn RSNA 2024. \emph{Tuy nhiên do giới hạn thời gian và GPU, báo cáo chỉ thực nghiệm và đánh giá trên 3 conditions (Spinal Canal Stenosis + L/R Foraminal Narrowing) --- 2 nhãn Subarticular được để cho phiên bản camera-ready.}
        \item \textbf{Trên RSNA 2024} (single-seed, full val), Hybrid cải thiện rõ rệt so với Baseline ResNet-34 thuần: Mean F1 macro tăng từ 0.343 lên \textbf{0.474} (+0.131 absolute), Mean AUC từ 0.766 lên 0.818, Mean AUPRC từ 0.530 lên 0.601.
        \item Cải thiện đáng kể nhất nằm ở \textbf{Severe class}: Severe Recall tăng từ 11.5\% (Baseline) lên \textbf{46.4\%} (Hybrid) --- gấp 4 lần. Severe F1 tăng từ 0.149 lên 0.343. Đây là cải tiến có ý nghĩa lâm sàng rõ rệt vì Severe chính là nhóm bệnh nặng cần ưu tiên phát hiện.
        \item Trên SPIDER 8-label transfer learning (3-seed evaluation), kết quả phát hiện một finding quan trọng: \textbf{CBAM-only hoạt động ngược cross-dataset} (Mean F1 0.619 $\pm$ 0.010 thấp hơn Baseline 0.646 $\pm$ 0.002), trong khi \textbf{Hybrid cân bằng được tốt hơn} (Mean F1 0.653 $\pm$ 0.014). BiomedCLIP frozen branch đóng vai trò regularizer cho cross-dataset generalization.
        \item Đóng góp BiomedCLIP ổn định ở các minority class: Disc Herniation Recall trên SPIDER tăng từ 0.298 (Baseline) lên 0.356 (Hybrid), trong đó BMC-only đã đóng góp 0.270 --- chứng minh prior từ vision-language model có giá trị cho lớp hiếm.
        \item Pipeline xử lý mất cân bằng lớp (FocalLoss + UncertaintyLoss + minority oversampling) đã được triển khai và chứng minh hiệu quả trong việc cải thiện Severe Recall.
    \end{itemize}
\section{Đề xuất hướng phát triển}
Để khắc phục các hạn chế hiện tại và hoàn thiện hệ thống hướng tới ứng dụng thực tế, hai hướng nghiên cứu trọng tâm tiếp theo được đề xuất:

\subsection{Hoàn thiện đánh giá đầy đủ 5/5 conditions của RSNA}
Kiến trúc Hybrid đề xuất đã được thiết kế với 5 classification head bao gồm cả Left/Right Subarticular Stenosis (2 nhãn mới của RSNA 2024) --- đây là phần mà SpineNetV2 nguyên bản còn thiếu. Tuy nhiên, do giới hạn thời gian và GPU, báo cáo chỉ thực nghiệm và đánh giá trên 3 conditions (Spinal Canal Stenosis + L/R Foraminal Narrowing).

Hướng phát triển ngắn hạn: huấn luyện và đánh giá Hybrid model trên đủ 5/5 conditions, bổ sung thêm nhánh xử lý đặc trưng từ mặt cắt Axial (nơi quan sát Subarticular rõ nhất) để đạt độ bao phủ bệnh lý toàn diện theo chuẩn RSNA 2024.

\subsection{Tối ưu hóa TotalSpineSeg cho mặt cắt Axial}
Mặc dù TotalSpineSeg đang hoạt động tốt trong việc định vị thân đốt sống trên mặt phẳng dọc, nhưng chẩn đoán các bệnh lý như hẹp ống sống hay hẹp ngách bên lại phụ thuộc rất lớn vào thông tin từ mặt cắt ngang .

Do đó, cần nghiên cứu huấn luyện lại hoặc tinh chỉnh  TotalSpineSeg để thực hiện phân vùng chính xác các cấu trúc (đĩa đệm, ống sống, rễ thần kinh) trên ảnh Axial. Việc có được mặt nạ phân vùng chính xác trên cả hai mặt phẳng sẽ giúp hệ thống định vị vùng tổn thương (ROI) chính xác hơn gấp nhiều lần, từ đó nâng cao đáng kể độ chính xác của mô hình chẩn đoán bệnh lý.

\subsection{Hoàn thiện phần mềm với giao diện cho bác sĩ}
Một hướng phát triển trực tiếp là hoàn thiện hệ thống thành một sản phẩm phần mềm cho bác sĩ. Các tính năng chính bao gồm: (a) giao diện trực quan để xem kết quả phân vùng và highlight vùng bất thường ngay trên ảnh DICOM, (b) cho phép bác sĩ chỉnh sửa nhãn và điều chỉnh vùng khoanh, (c) xuất ảnh đã gán nhãn dưới dạng độc lập, thuận tiện cho việc lưu trữ và chia sẻ. Sản phẩm sẽ tích hợp pipeline 2 mô-đun đã xây dựng (TotalSpineSeg + Hybrid grading) thành một quy trình end-to-end. Dự kiến hoàn thiện trong giai đoạn nghiên cứu tiếp theo.

\subsection{Đánh giá đa-seed cho RSNA và bootstrap CI}
Báo cáo hiện tại chỉ chạy single-seed cho RSNA 2024 do hạn chế thời gian GPU. Hướng phát triển ngắn hạn là mở rộng 3-seed evaluation (seeds 42, 123, 456) trên RSNA tương tự cách đã làm cho SPIDER, đồng thời bổ sung bootstrap 95\% CI ($n = 1000$ resamples) trên Severe F1, Severe Recall và Severe AUPRC. Việc này sẽ tăng độ tin cậy của các chỉ số headline (Severe Recall 11.5\% $\rightarrow$ 46.4\%) bằng cách định lượng được mức độ biến thiên.

\subsection{Tích hợp Vision-Language Model cho việc sinh báo cáo chẩn đoán}
Một hướng phát triển khác là tích hợp một mô hình Vision-Language lớn (VLM) như MedGemma vào pipeline để tự động sinh báo cáo chẩn đoán bằng ngôn ngữ tự nhiên. Trong khảo sát của báo cáo (Chương 5), MedGemma nguyên bản có Recall rất thấp (6--21\%) trên bệnh lý thắt lưng do chưa được chuyên biệt hóa nên không được đưa vào pipeline đề xuất. Hai hướng cải thiện cho việc tích hợp VLM:
\begin{itemize}
    \item \textbf{Fine-tune trên dữ liệu MRI thắt lưng:} dùng các cặp (ảnh, văn bản chẩn đoán) từ RSNA và SPIDER để fine-tune nhánh thị giác và nhánh ngôn ngữ của VLM. Mục tiêu cải thiện độ nhạy với các thuật ngữ và pattern hình ảnh cụ thể của bệnh thắt lưng.
    \item \textbf{Tích hợp VLM với output phân vùng:} mở rộng prompt cho VLM để bao gồm cả mặt nạ phân vùng từ TotalSpineSeg (overlay màu lên ảnh gốc), giúp VLM có thông tin spatial rõ ràng hơn khi viết báo cáo chẩn đoán. Hướng này cũng có thể được mở rộng để xử lý cả mặt cắt Axial (ngang) bên cạnh mặt cắt Sagittal hiện tại.
\end{itemize}

\subsection{Cải thiện qua phản hồi từ bác sĩ chẩn đoán (HITL)}
Một hướng cải tiến dài hạn là tích hợp cơ chế \textbf{Human-in-the-Loop} (HITL). Cụ thể: thông qua giao diện phần mềm tương lai, bác sĩ có thể chỉnh sửa các nhãn AI gán sai trong quá trình sử dụng thực tế. Các cặp ảnh + nhãn đã được chỉnh sửa sẽ được lưu lại dưới dạng dataset mới, dùng để huấn luyện lại (retrain) mô hình theo định kì, tạo ra một vòng lặp cải tiến liên tục. Hướng này hiện chưa nằm trong scope LVTN do đòi hỏi sự phối hợp với bệnh viện thực tế và một khoảng thời gian đủ dài để thu thập đủ phản hồi --- được đặt như định hướng nghiên cứu cho các giai đoạn tiếp theo.